\scanpage{132}

% L1
\textcircled{4} \uline{Cas $n$ pair} $=2\nu$

% L2
\[
 q_n=q_{2\nu}=\sum_1^\nu(x_i y_i)\qquad \delta'(q_n)=(-1)^\nu\tag*{(4.1)}
\]

% L3
\deleted{\unsure{(4|(4.}} Ici on a

% L4–6
\[
 \underset{\substack{\parallel\\\mathrm{Centre}\,O(q)}}{\mu_2}
 \subset\underbrace{\deleted{\unsure{S}}O(q)\cap\mathrm{Sl}(E)}_{
 \substack{\text{on ne l'appelle}\\\text{\uline{pas}}\ SO(q)\\\text{à cause de car. }2}}
 \tag*{(4.2)}
\]

% L7
\deleted{et \unsure{on a}}

% L8–9, left rule
\uline{Thm} Pour $n$ pair, $O(q)$ est lisse, si $SO(q)$ est sa composante neutre, on a

% L10–11
\[
 \begin{aligned}
 &O(q)/SO(q)\simeq(\mathbb Z/2)_S\\
 \text{et}\qquad &\mu_2=\mathrm{Centre}(O(q))\subset SO(q).
 \end{aligned}
\]

% L12–13
\uline{Cor} Si $X=V(q)$, le groupe $\underline{\mathrm{Aut}}(X)\simeq\underline O(q)/\mu_2$ est lisse, et

% L14–15
\[
 \left\{\begin{aligned}
 &\underline{\mathrm{Aut}}(X)^{\unsure{0|\circ}}\simeq SO(q)/\mu_2\\
 &\underline{\mathrm{Aut}}(X)/\underline{\mathrm{Aut}}(X)^0\simeq(\mathbb Z/2)_S
 \end{aligned}\right.
\]

% L16
\unsure{Indication|Indications} de dém. du Thm.

% L17–18
La chose essentielle est de définir l'hom \deleted{\unsure{surj}} épi

% L18
\[
 O(q)\longrightarrow(\mathbb Z/2)_S
\]

% L19–20
et pour ceci on va associer can. à $q$ un rev. étale de degré $2$ de $S$.

% lower left
% archive sheet number: 131
